\subsection{Economic interpretation and the NBO program}
The new evidence connects a feasible implementation to two distinct economic statements. First, the acquisition-aware theorem bounds the loss of a policy that can actually be executed under uncertain state measurement and capacity constraints. Second, the direct experiment certifies proximity of the original scalar policy costs within an explicitly chosen decision tolerance, so that a user of that catalogue may consider implementation charges without treating separate regret bounds as values. These statements are stronger than a certificate comparison alone.

Neither statement implies that a neural representation is indispensable. The identical native envelope is a mathematical control, and the conventional fitted-value and spline implementations remain strong alternatives. The constrained experiment demonstrates two-state continuous-uncertainty construction and nontrivial repair, not a high-dimensional scaling result or an estimated economic application. Its tight targets remain unachieved. The direct cost experiment does not cover the new constrained policies or nonlinear recursive utility; its expectation-specific conclusion is not transferred to those settings. The controlled-diffusion, recursive-preference, endogenous-preference, temporal-self and game developments below retain their own conditions and original role in the NBO program. Their preservation is not counted as a new comparative execution.
